 %%%%%%%%%%%%%%%%%%%%%%%%%%%%%%%%%%%%%%%%%%%%%%%%%%%
\section{Тепловые процессы}
%%%%%%%%%%%%%%%%%%%%%%%%%%%%%%%%%%%%%%%%%%%%%%%%%%%%
\subsection{Тепловое расширение}
%%%%%%%%%%%%%%%%%%%%%%%%%%%%%%%%%%%%%%%%%%%%%%%%%%%%
%%%%%%%%%%%%%%%%%%%%%%%%%%%%%%%%%%%%%%%%%%%%%%%%%%%%
\z{При температуре $t_1 = 20\degree\text{C}$ длины алюминиевого и медного стержней одинаковы. Какова длина алюминиевого стержня при $t_2 = -20\degree\text{C}$, если длина медного стержня при этой температуре $l_{\text{Cu}} = 60\,\text{см}$? Коэффициенты линейного теплового расширения: алюминия -- $\alpha_{\text{Al}} = 23.8 \cdot 10^{-6}\,\text{К}^{-1}$; меди -- $\alpha_{\text{Cu}} = 16.5 \cdot 10^{-6}\,\text{К}^{-1}$.}
%
{$l_{Cu2}=59.98 \,\text{см}$}

%%%%%%%%%%%%%%%%%%%%%%%%%%%%%%%%%%%%%%%%%%%%%%%%%%%%
\z{Из медного диска радиуса $R$ вырезали круг радиуса $r$. Диск нагревают от $T_1$ до $T_2$. Какой станет площадь дырочки?}
%
{$S_2 = \pi r_0^2 (1 + \alpha \Delta T)^2$}

%%%%%%%%%%%%%%%%%%%%%%%%%%%%%%%%%%%%%%%%%%%%%%%%%%%%
 \z{Как должны относиться длины $L_1$ и $L_2$ двух стержней, сделанных из разных материалов, с температурными коэффициентами $\alpha_1$ и $\alpha_2$, чтобы при любой температуре разность длин стержней оставалась одинаковой?}
%
{$\frac{L_1}{L_2} = \frac{\alpha_2}{\alpha_1}$}

%%%%%%%%%%%%%%%%%%%%%%%%%%%%%%%%%%%%%%%%%%%%%%%%%%%%
 \z{Диаметр стеклянной пробки, застрявшей в горлышке флакона, $d_0 = 2.5\,\text{см}$. Чтобы вынуть пробку, горлышко нагрели до температуры $t_1 = 150\degree\text{С}$. Сама пробка успела при этом нагреться до температуры $t_2 = 50\degree\text{С}$. Как велик образовавшийся зазор? Температурный коэффициент линейного расширения стекла $\alpha_1 = 9 \cdot 10^{-6}\,\text{К}^{-1}$.}
%
{$l=\frac{d_0\alpha_1(t_1-t_2)}{2}=0.01\,\text{мм}$}

%%%%%%%%%%%%%%%%%%%%%%%%%%%%%%%%%%%%%%%%%%%%%%%%%%%%
\z{Плотность ртути уменьшилась при нагревании до $98\%$ от её плотности при $0\,\degree\text{C}$. До какой температуры нагрели ртуть? Коэффициент теплового объёмного расширения ртути $\beta_{\text{Hg}} = 0.181 \cdot 10^{-3}\,\text{К}^{-1}$.}
%
{$T_{кон}=110.5\,\degree\text{C}$}

%%%%%%%%%%%%%%%%%%%%%%%%%%%%%%%%%%%%%%%%%%%%%%%%%%%%
\z{При температуре $t_0 = 0\,\degree\text{С}$ стеклянный баллон вмещает $m_0 = 100\,\text{г}$ ртути. При температуре $t_1 = 20\,\degree\text{С}$ баллон вмещает $m_1 = 99.7\,\text{г}$ ртути. В обоих случаях температура ртути равна температуре баллона. Найдите по этим данным температурный коэффициент линейного расширения стекла $\alpha_1$, учитывая, что коэффициент объемного расширения ртути $\alpha = 1.8 \cdot 10^{-4}\,\text{К}^{-1}$.}
%
{$\alpha_1=\frac{m_1(1+\alpha t_1)-m_0}{m_0t_1}=10^{-5}\,\text{К}^{-1}$}
	
%%%%%%%%%%%%%%%%%%%%%%%%%%%%%%%%%%%%%%%%%%%%%%%%%%%%
\z{Толщина биметаллической пластинки, составленной из одинаковых полосок стали и цинка, равна $d = 0.1\,\text{см}$. Определите радиус кривизны $R$ пластинки при повышении температуры на $\Delta t = 11\,\degree\text{C}$. Коэффициент линейного расширения цинка $\alpha_1 = 25\cdot10^{-6}\,\text{К}^{-1}$, стали $\alpha_2 = 12\cdot10^{-6}\,\text{К}^{-1}$.}
%
{$R \approx 350\ \text{см}$}




